\documentclass[a4paper,12pt]{article}
\usepackage[utf8]{inputenc}
\usepackage[russian]{babel}
\usepackage{graphicx}
\usepackage{listings}
\usepackage{hyperref}
\usepackage{geometry}
\geometry{top=2cm, bottom=2cm, left=3cm, right=1.5cm}

\title{Отчет по практической работе №1\\Дисциплина: GVS}
\author{Студент}
\date{\today}

\begin{document}
\maketitle

\section{Цель работы}
Освоить базовые навыки программирования CUDA: работу с одномерными сетками нитей и динамической памятью устройства. Изучить паттерн проектирования data + view.

\section{Фрагменты исходного кода}
\subsection{Data}
\begin{lstlisting}[language=C++, basicstyle=\ttfamily\scriptsize]
// core/include/Data.cuh
template <typename AtomT>
class Data { ... };
\end{lstlisting}

\subsection{VectorView}
\begin{lstlisting}[language=C++, basicstyle=\ttfamily\scriptsize]
// core/include/VectorView.cuh
template <typename AtomT>
class VectorView { ... };
\end{lstlisting}

\subsection{Vector}
\begin{lstlisting}[language=C++, basicstyle=\ttfamily\scriptsize]
// core/include/Vector.cuh
template <typename AtomT>
class Vector { ... };
\end{lstlisting}

\subsection{kernel\_vecadd}
\begin{lstlisting}[language=C++, basicstyle=\ttfamily\scriptsize]
// core/include/kernel_vecadd.cuh
template <typename AtomT>
__global__ void kernel_vecadd(VectorView<AtomT> lhs, 
                              VectorView<AtomT> rhs, 
                              VectorView<AtomT> out) { ... }
\end{lstlisting}

\section{Ссылка на репозиторий}
\href{https://github.com/your-username/gvs}{Репозиторий с полным кодом}

\section{Результаты измерений}
% Вставьте сюда сгенерированные графики:
% \begin{figure}[h]
% \centering
% \includegraphics[width=0.8\textwidth]{real_complexity.png}
% \caption{График реальной вычислительной сложности}
% \end{figure}

% \begin{figure}[h]
% \centering
% \includegraphics[width=0.8\textwidth]{speedup.png}
% \caption{График ускорения}
% \end{figure}

\section{Интерпретация результатов}
Теоретическая сложность сложения векторов линейна $O(N)$. 
График реальной сложности подтверждает это: для CPU (Eigen) время растет пропорционально количеству элементов.
Для GPU (CUDA) наблюдается накладной расход на запуск ядра при малых $N$, поэтому время почти константно. Однако при больших $N$ рост также становится линейным, но с меньшим коэффициентом пропорциональности.
График ускорения показывает, что на малых $N$ CPU выигрывает из-за накладных расходов CUDA API. Начиная с некоторого размера, GPU начинает показывать значительное ускорение, выходящее на плато (ограничено пропускной способностью памяти).

\end{document}

